% GENERATED FILE. Edit canonical lessons, then run npm run generate:books.
% canonical-source-hash: fnv1a64:692c3eaa2c470081
% canonical-lessons: GU-C551-talim, GU-C551-koc, GU-C551-riyajh, GU-C551-dodavvu, GU-C551-trophi

\chapter{Training for the Big Day}
\label{ch:gu-training-for-the-big-day}
\hlchaptermodality{551}

\begin{chapteropening}
\textbf{By the end of this chapter:} \emph{I can talk about training and practice before a big day.}

Complete the last lesson of chapter 551: I can talk about training and practice before a big day.
\end{chapteropening}

\section[\texorpdfstring{\gu{તાલીમ} (tālīm)}{તાલીમ (tālīm)}]{\gu{તાલીમ} (tālīm) — training}

\label{lesson:GU-C551-talim}



\begin{quote}
Before the new one: say the Gujarati for to pierce (a target), then the Gujarati for a runner.
\end{quote}

\subsection*{You'll want to know: \gu{તાલીમ}}
\textbf{\gu{તાલીમ}} — \emph{tālīm} — \textquotedblleft{}training\textquotedblright{}.

\begin{quote}
\gu{જો} \gu{તમે} \gu{સારી} \gu{તાલીમ} \gu{લેશો} \gu{તો} \gu{તમે} \gu{ઝડપી} \gu{દોડી} \gu{શકશો}.
\end{quote}

\begin{quote}
\emph{jo tame sārī tālīm leśo to tame jhaḍpī doḍī śakśo.}
\end{quote}

\begin{quote}
If you get good training, you'll be able to run fast.
\end{quote}

A feminine noun.

\begin{cousinweb}
\textbf{\gu{તાલીમ}} (\emph{tālīm}) is usually traced to Arabic \emph{taʿlīm}, \textquotedblleft{}teaching, instruction\textquotedblright{}, through Persian.
\end{cousinweb}

\begin{grammarlens}[title={What you've built}]
One more word for saying what happens if.
\end{grammarlens}

\subsection*{Your turn}
\begin{itemize}
\raggedright
  \item \emph{Say it:} \emph{tālīm}
  \item \emph{Say it:} \emph{tālīm}, once more
  \item \emph{Say it:} \emph{jo tame sārī tālīm leśo to tame jhaḍpī doḍī śakśo.} — If you get good training, you'll be able to run fast.
  \item \emph{Recall:} say the Gujarati for to pierce (a target), then the Gujarati for a runner, then say \emph{tālīm} again
\end{itemize}

\begin{tcolorbox}[breakable,colback=teal!4,colframe=teal!35!black,title={Before you move on}]
What does \gu{તાલીમ} mean? (\textquotedblleft{}Training\textquotedblright{}.) Say it once more.
\end{tcolorbox}
\section[\texorpdfstring{\gu{કોચ} (koc)}{કોચ (koc)}]{\gu{કોચ} (koc) — a sports coach}

\label{lesson:GU-C551-koc}



\begin{quote}
Before the new one: say the Gujarati for a runner, then the Gujarati for training.
\end{quote}

\subsection*{You'll want to know: \gu{કોચ}}
\textbf{\gu{કોચ}} — \emph{koc} — \textquotedblleft{}a sports coach\textquotedblright{}.

\begin{quote}
\gu{જો} \gu{કોચ} \gu{ના} \gu{કહેશે} \gu{તો} \gu{હું} \gu{મેચમાં} \gu{નહીં} \gu{રમું}.
\end{quote}

\begin{quote}
\emph{jo koc nā kaheśe to hũ mecmā̃ nahī̃ ramũ.}
\end{quote}

\begin{quote}
If the coach says no, I won't play in the match.
\end{quote}

A masculine or feminine noun, by the person meant.

\begin{cousinweb}
\textbf{\gu{કોચ}} (\emph{koc}) is borrowed from English \emph{coach}.
\end{cousinweb}

\begin{grammarlens}[title={What you've built}]
One more word for saying what happens if.
\end{grammarlens}

\subsection*{Your turn}
\begin{itemize}
\raggedright
  \item \emph{Say it:} \emph{koc}
  \item \emph{Say it:} \emph{koc}, once more
  \item \emph{Say it:} \emph{jo koc nā kaheśe to hũ mecmā̃ nahī̃ ramũ.} — If the coach says no, I won't play in the match.
  \item \emph{Recall:} say the Gujarati for a runner, then the Gujarati for training, then say \emph{koc} again
\end{itemize}

\begin{tcolorbox}[breakable,colback=teal!4,colframe=teal!35!black,title={Before you move on}]
What does \gu{કોચ} mean? (\textquotedblleft{}A sports coach\textquotedblright{}.) Say it once more.
\end{tcolorbox}
\section[\texorpdfstring{\gu{રિયાઝ} (riyājh)}{રિયાઝ (riyājh)}]{\gu{રિયાઝ} (riyājh) — music practice}

\label{lesson:GU-C551-riyajh}



\begin{quote}
Before the new one: say the Gujarati for training, then the Gujarati for a sports coach.
\end{quote}

\subsection*{You'll want to know: \gu{રિયાઝ}}
\textbf{\gu{રિયાઝ}} — \emph{riyājh} — \textquotedblleft{}music practice\textquotedblright{}.

\begin{quote}
\gu{જો} \gu{મીરા} \gu{રોજ} \gu{રિયાઝ} \gu{કરશે} \gu{તો} \gu{તેનો} \gu{સૂર} \gu{વધારે} \gu{મીઠો} \gu{થશે}.
\end{quote}

\begin{quote}
\emph{jo Mīrā roj riyājh karśe to teno sūr vadhāre mīṭho thaśe.}
\end{quote}

\begin{quote}
If Mira practises every day, her singing will get sweeter.
\end{quote}

A masculine noun.

\begin{cousinweb}
\textbf{\gu{રિયાઝ}} (\emph{riyājh}) is usually traced to Arabic \emph{riyāḍa}, \textquotedblleft{}exercise, training\textquotedblright{}, through Persian \emph{riyāz}; a musician does it: \textbf{\gu{રિયાઝ} \gu{કરવો}}.
\end{cousinweb}

\begin{grammarlens}[title={What you've built}]
One more word for saying what happens if.
\end{grammarlens}

\subsection*{Your turn}
\begin{itemize}
\raggedright
  \item \emph{Say it:} \emph{riyājh}
  \item \emph{Say it:} \emph{riyājh}, once more
  \item \emph{Say it:} \emph{jo Mīrā roj riyājh karśe to teno sūr vadhāre mīṭho thaśe.} — If Mira practises every day, her singing will get sweeter.
  \item \emph{Recall:} say the Gujarati for training, then the Gujarati for a sports coach, then say \emph{riyājh} again
\end{itemize}

\begin{tcolorbox}[breakable,colback=teal!4,colframe=teal!35!black,title={Before you move on}]
What does \gu{રિયાઝ} mean? (\textquotedblleft{}Music practice\textquotedblright{}.) Say it once more.
\end{tcolorbox}
\section[\texorpdfstring{\gu{દોડાવવું} (doḍāvvũ)}{દોડાવવું (doḍāvvũ)}]{\gu{દોડાવવું} (doḍāvvũ) — to make (someone) run}

\label{lesson:GU-C551-dodavvu}



\begin{quote}
Before the new one: say the Gujarati for a sports coach, then the Gujarati for music practice.
\end{quote}

\subsection*{You'll want to know: \gu{દોડાવવું}}
\textbf{\gu{દોડાવવું}} — \emph{doḍāvvũ} — \textquotedblleft{}to make (someone) run\textquotedblright{}.

\begin{quote}
\gu{જો} \gu{કોચ} \gu{અમને} \gu{રોજ} \gu{ન} \gu{દોડાવત} \gu{તો} \gu{અમે} \gu{આટલા} \gu{ઝડપી} \gu{ન} \gu{હોત}.
\end{quote}

\begin{quote}
\emph{jo koc amne roj na doḍāvat to ame āṭlā jhaḍpī na hot.}
\end{quote}

\begin{quote}
If the coach hadn't made us run every day, we wouldn't be this fast.
\end{quote}

Past, with \gu{તેણે} (he or she …): \gu{તેણે} \gu{છોકરાઓને} \gu{દોડાવ્યા}, \emph{teṇe chokrāone doḍāvyā}.

\begin{cousinweb}
\textbf{\gu{દોડાવવું}} (\emph{doḍāvvũ}) is built from \textbf{\gu{દોડવું}}, \textquotedblleft{}to run\textquotedblright{}, with the causative \emph{-āv-}. CDIAL files \textbf{\gu{દોડવું}} under Sanskrit \emph{drávati}, \textquotedblleft{}runs\textquotedblright{} (CDIAL 6624).
\end{cousinweb}

\begin{grammarlens}[title={What you've built}]
One more word for saying what happens if.
\end{grammarlens}

\subsection*{Your turn}
\begin{itemize}
\raggedright
  \item \emph{Say it:} \emph{doḍāvvũ}
  \item \emph{Say it:} \emph{doḍāvvũ}, once more
  \item \emph{Say it:} \emph{jo koc amne roj na doḍāvat to ame āṭlā jhaḍpī na hot.} — If the coach hadn't made us run every day, we wouldn't be this fast.
  \item \emph{Recall:} say the Gujarati for a sports coach, then the Gujarati for music practice, then say \emph{doḍāvvũ} again
\end{itemize}

\begin{tcolorbox}[breakable,colback=teal!4,colframe=teal!35!black,title={Before you move on}]
What does \gu{દોડાવવું} mean? (\textquotedblleft{}To make (someone) run\textquotedblright{}.) Say it once more.
\end{tcolorbox}
\section[\texorpdfstring{\gu{ટ્રોફી} (ṭrophī)}{ટ્રોફી (ṭrophī)}]{\gu{ટ્રોફી} (ṭrophī) — a trophy}

\label{lesson:GU-C551-trophi}



\begin{quote}
Before the new one: say the Gujarati for music practice, then the Gujarati for to make (someone) run.
\end{quote}

\subsection*{You'll want to know: \gu{ટ્રોફી}}
\textbf{\gu{ટ્રોફી}} — \emph{ṭrophī} — \textquotedblleft{}a trophy\textquotedblright{}.

\begin{quote}
\gu{જો} \gu{અમારી} \gu{ટીમ} \gu{જીતશે} \gu{તો} \gu{શાળાને} \gu{ટ્રોફી} \gu{મળશે}.
\end{quote}

\begin{quote}
\emph{jo amārī ṭīm jītśe to śāḷāne ṭrophī maḷśe.}
\end{quote}

\begin{quote}
If our team wins, the school will get a trophy.
\end{quote}

A feminine noun.

\begin{cousinweb}
\textbf{\gu{ટ્રોફી}} (\emph{ṭrophī}) is borrowed from English \emph{trophy}.
\end{cousinweb}

\begin{grammarlens}[title={What you've built}]
One more word for saying what happens if.
\end{grammarlens}

\subsection*{Your turn}
\begin{itemize}
\raggedright
  \item \emph{Say it:} \emph{ṭrophī}
  \item \emph{Say it:} \emph{ṭrophī}, once more
  \item \emph{Say it:} \emph{jo amārī ṭīm jītśe to śāḷāne ṭrophī maḷśe.} — If our team wins, the school will get a trophy.
  \item \emph{Recall:} say the Gujarati for music practice, then the Gujarati for to make (someone) run, then say \emph{ṭrophī} again
\end{itemize}

\begin{tcolorbox}[breakable,colback=teal!4,colframe=teal!35!black,title={Before you move on}]
What does \gu{ટ્રોફી} mean? (\textquotedblleft{}A trophy\textquotedblright{}.) Say it once more.
\end{tcolorbox}
